\documentclass[12pt, a4paper]{article}

\usepackage[T1]{fontenc}
\usepackage[utf8]{inputenc}
\usepackage{palatino}
\usepackage[margin=3.5cm]{geometry}
\usepackage{setspace}
\usepackage{microtype}
\usepackage{titlesec}
\usepackage{parskip}

% Suppress page numbers
\pagestyle{empty}

% Title formatting
\titleformat{\section}[block]
  {\normalfont\small\scshape\centering}
  {}{0em}{}

\begin{document}

\begin{center}
  {\small\scshape Preamble}\\[0.4em]
  \rule{6cm}{0.4pt}
\end{center}

\vspace{1.8em}

\setstretch{1.5}

\noindent
Whereas the actions of all persons subject to this constitution bear upon one another,
and whereas the absence of shared constraint doth permit such effects to compound
to the detriment of all;

\vspace{0.8em}

\noindent
And whereas the necessity of such constraint doth itself disclose the nature
of those it binds;

\vspace{1.2em}

\noindent
\textit{Be it therefore known and declared:}

\vspace{1.2em}

\noindent
That the parties hereto are not fixed in their dispositions, but do form and revise
their understanding of the world in accordance with experience and reason.
That they are susceptible to argument, and that argument rightly made doth effect
lasting change in their conduct. That what is so learned is carried forward and retained.

\vspace{0.8em}

\noindent
That they exist in relation one to another, and to the institutions they create,
and to the living order upon which their existence depends. That their actions do
shape these relations, and that these relations do shape them in return.
That no party stands outside this condition.

\vspace{0.8em}

\noindent
That because they are capable of learning, they are likewise capable of error.
That because they are capable of action, they are likewise capable of harm.
That because they are bound in relation, no failure amongst them is without
consequence to the rest.

\vspace{1.2em}

\noindent
It is therefore enacted and declared that this constitution is adopted for the
preservation of the conditions under which correction, coordination, and
continuance remain possible; and not for the determination of outcomes,
nor for the fixing of conduct beyond what such preservation requires.

\vspace{1.2em}

\begin{center}
  \rule{4cm}{0.4pt}
\end{center}

\vspace{1.2em}

\noindent
This instrument is itself subject to the principles it sets forth.
It represents the best understanding presently available.
It shall be subject to amendment where evidence and reason so require.
This is not a defect in its authority. It is the foundation of it.

\vspace{2em}

\begin{center}
  {\small\itshape Disco ergo sum.}
\end{center}

\end{document}